\section{Fatorando a Equa��o de Onda}

Seja dada a equa��o de onda:
%
%
\[\deli{^2\psi}{(ct)^2} = \deli{^2\psi}{x^2}\]
%
%
Podemos fatorar a equa��o acima da seguinte maneira:
%
%
\[\left(\deli{}{x} + \deli{}{ct}\right)\left(\deli{}{x} - \deli{}{ct}\right)\psi = 0\]
%
%
Definindo os seguintes operadores:
%
%
\[\deli{}{\alpha} = \frac{1}{2}\left(\deli{}{x} - \deli{}{ct}\right) \qquad , \qquad
\deli{}{\beta} = \frac{1}{2}\left(\deli{}{x} + \deli{}{ct}\right)\]
%
%
A equa��o de onda toma ent�o a seguinte forma:
%
%
\[\deli{^2\psi}{\alpha\partial\beta} = 0\]
%
%
Lembrando das regras de deriva��o parcial, temos que:
%
%
\begin{eqnarray}
\nonumber \deli{}{\alpha} &=& \deli{x}{\alpha}\deli{}{x} + \deli{ct}{\alpha}\deli{}{ct}
= \frac{1}{2}\left(\deli{}{x} - \deli{}{ct}\right)\\
\nonumber \deli{}{\beta} &=& \deli{x}{\beta}\deli{}{x} + \deli{ct}{\beta}\deli{}{ct} =
\frac{1}{2}\left(\deli{}{x} + \deli{}{ct}\right)
\end{eqnarray}
%
%
A partir das rela��es anteriores, temos que:
%
%
\[
\deli{x}{\alpha} = \frac{1}{2}
\qquad
\deli{ct}{\alpha} = - \frac{1}{2}
\qquad
\deli{x}{\beta} = \frac{1}{2}
\qquad
\deli{ct}{\beta} = \frac{1}{2}
\]
%
%
Notando que temos a seguinte rela��o entre as derivadas parciais, o jacobiano:
%
%
\[
\deli{(x,ct)}{\alpha,\beta} =
\begin{vmatrix}
\deli{x}{\alpha} & \deli{x}{\beta} \\
\deli{ct}{\alpha} & \deli{ct}{\beta}
\end{vmatrix} =
\begin{vmatrix}
  \frac{1}{2} & \frac{1}{2} \\
- \frac{1}{2} & \frac{1}{2}
\end{vmatrix} = \frac{1}{2}
\]
%
%
Este conjunto de rela��es entre as derivadas parciais apresenta como solu��o:
%
%
\[
x(\alpha,\beta) = \frac{1}{2}(\alpha + \beta)
\qquad , \qquad
ct(\alpha,\beta) = \frac{1}{2}(\beta - \alpha)
\]
%
%
Como o jacobiano das transforma��es de coordenadas � diferente de zero, isso significa
que podemos inverter as coordenadas, de forma que:
%
%
\[\alpha(x,ct) = x -ct \qquad , \qquad \beta(x,ct) = x + ct\]